%% FullCourse_10Lecs -- Lecture 9, Topic 9.3: Homotopy Workflow + Tool Choice
%% Source: ./archive_pre_split/Lec09_Agentic_AI_v3.tex (sections 6-7)
%% Pass B: layer in refinements from
%%         ../../MiniCourse_8hr/Slides/Module4_Agentic_AI/M4_T9_Homotopy_Failures_Safety.tex
%% Rewritten 2026-07-06: Act 3 "Method" of the five-act structure. Six-step pipeline
%%         labels kept VERBATIM (Lec10 T8 reproduces them); adds HA-Ramsey opening
%%         prompts (Lec10 T8 F18), seed criteria (F12), three-layers frame (F11),
%%         and the "Grade the Plan" criteria frame. See Lec09_Rewrite_Plan_2026-07-06.md.

%% Shared preamble for the FullCourse_10Lecs topic decks.
%%
%% Each topic .tex \input's this file, then defines its own \title, \date
%% override (if any), \graphicspath, and \begin{document}.
%%
%% Reconstructed verbatim from the inlined copy preserved in
%% Lec10_Case_Studies/Lec10_T8_Case6_DeepRamsey.tex (lines 23-190), which is
%% explicitly annotated there as "from ../shared_preamble.tex, copied verbatim".
%% Base preamble matches MiniCourse_8hr/Slides/shared_preamble.tex; this version
%% additionally provides \sectiondivider, the conceptbox/cautionbox/examplebox/
%% takeawaybox tcolorboxes, and \compacttables used across the split decks.

\documentclass[10pt,english,aspectratio=169]{beamer}

\usepackage{ae,aecompl}
\usepackage[T1]{fontenc}
\usepackage{textcomp}
\usepackage[utf8]{inputenc}
\usepackage{babel}
\usepackage{datetime2}

\setcounter{secnumdepth}{3}
\setcounter{tocdepth}{3}

\usepackage{amsmath, amssymb, amsthm, graphicx}
\usepackage{booktabs, multirow}
\usepackage{tikz}
\usepackage{pgfplots}
\pgfplotsset{compat=1.16}
\usepackage[most]{tcolorbox}
\tcbuselibrary{skins, breakable, theorems}
\usepackage{listings}
\usepackage{xcolor}

\lstset{
  basicstyle=\ttfamily\small,
  keywordstyle=\color{blue!80!black}\bfseries,
  commentstyle=\color{green!50!black}\itshape,
  stringstyle=\color{orange!80!black},
  backgroundcolor=\color{gray!8},
  frame=single,
  rulecolor=\color{gray!40},
  breaklines=true,
  showstringspaces=false,
  numbers=none,
}

\lstdefinestyle{pythonstyle}{
    language=Python,
    basicstyle=\ttfamily\footnotesize,
    keywordstyle=\color{blue!80!black}\bfseries,
    commentstyle=\color{green!50!black}\itshape,
    stringstyle=\color{orange!80!black},
    frame=single,
    breaklines=true,
    showstringspaces=false,
    numbers=left,
    numbersep=5pt,
    numberstyle=\tiny\color{gray},
    backgroundcolor=\color{gray!8},
    rulecolor=\color{gray!40}
}

\ifx\hypersetup\undefined
  \AtBeginDocument{%
    \hypersetup{bookmarks=false,breaklinks=true,pdfborder={0 0 0},colorlinks=true,
      linkcolor=DarkText,urlcolor=RedTitle}%
  }
\else
  \hypersetup{bookmarks=false,breaklinks=true,pdfborder={0 0 0},colorlinks=true,
    linkcolor=DarkText,urlcolor=RedTitle}%
\fi

\makeatletter
\newcommand\makebeamertitle{%
  \begin{frame}[plain]
    \vspace*{1.0cm}
    \centering
    {\usebeamerfont{title}\usebeamercolor[fg]{title}\inserttitle\par}
    \vspace{1.0cm}
    {\usebeamerfont{author}\usebeamercolor[fg]{author}\insertauthor\par}
    \vspace{0.2cm}
    {\usebeamerfont{date}\usebeamercolor[fg]{date}\insertdate\par}
  \end{frame}}%
\AtBeginDocument{%
  \let\origtableofcontents=\tableofcontents
  \def\tableofcontents{\@ifnextchar[{\origtableofcontents}{\gobbletableofcontents}}
  \def\gobbletableofcontents#1{\origtableofcontents}%
}

\usetheme{metropolis}
\usefonttheme{professionalfonts}

\definecolor{LightGray}{RGB}{230,230,230}
\definecolor{RedTitle}{RGB}{0,0,200}
\definecolor{VeryLightGray}{RGB}{245,245,245}
\definecolor{DarkText}{RGB}{0,0,0}

\setbeamercolor{background canvas}{bg=white}
\setbeamercolor{title}{fg=RedTitle}
\setbeamercolor{author}{fg=DarkText}
\setbeamercolor{date}{fg=DarkText}
\setbeamercolor{frametitle}{bg=VeryLightGray,fg=DarkText}
\setbeamercolor{normal text}{fg=DarkText}
\setbeamerfont{title}{size=\large}

\setbeamersize{text margin left=5mm, text margin right=5mm}
\addtobeamertemplate{frametitle}{}{\vspace{-0.5em}}
\newcommand{\reducevspace}{\vspace{-0.5cm}}

% Shared section divider and box styles for split topic decks.
\newcommand{\sectiondivider}[2]{%
\begin{frame}[plain,noframenumbering]
\begin{center}
\vfill
{\Large\textbf{#1}\par}
\vspace{0.35cm}
{\normalsize #2\par}
\vfill
\end{center}
\end{frame}
}

\newtcolorbox{conceptbox}[1][]{
  colback=blue!3,
  colframe=RedTitle!55,
  boxrule=0.35mm,
  arc=1mm,
  left=2mm,
  right=2mm,
  top=1mm,
  bottom=1mm,
  #1
}
\newtcolorbox{cautionbox}[1][]{
  colback=red!3,
  colframe=red!50!black,
  boxrule=0.35mm,
  arc=1mm,
  left=2mm,
  right=2mm,
  top=1mm,
  bottom=1mm,
  #1
}
\newtcolorbox{examplebox}[1][]{
  colback=green!3,
  colframe=green!45!black,
  boxrule=0.35mm,
  arc=1mm,
  left=2mm,
  right=2mm,
  top=1mm,
  bottom=1mm,
  #1
}
\newtcolorbox{takeawaybox}[1][]{
  colback=VeryLightGray,
  colframe=RedTitle!55,
  boxrule=0.35mm,
  arc=1mm,
  left=2mm,
  right=2mm,
  top=1mm,
  bottom=1mm,
  #1
}
\newcommand{\compacttables}{\renewcommand{\arraystretch}{1.15}}

\setbeamertemplate{navigation symbols}{}
\setbeamertemplate{footline}{%
  \hfill%
  \usebeamercolor[fg]{page number in head/foot}%
  \usebeamerfont{page number in head/foot}%
  \insertframenumber\,/\,\inserttotalframenumber\kern1em\vskip2pt%
}

\makeatother

\author{Zhigang Feng}
% "date with time": compile timestamp so each build is identifiable.
\date{\today\ \DTMcurrenttime}


% Suppress redundant auto-generated metropolis section page; \sectiondivider frames are the dividers we keep.
\metroset{sectionpage=none}

\graphicspath{{./pic/}{./figures/}{../../MiniCourse_8hr/Slides/Module4_Agentic_AI/pic/}{../}}


\usetikzlibrary{positioning,arrows.meta,shapes,calc}

%% Lec09 shared MAP slide: "one map, five acts" recurring diagram.
%% Input this file AFTER ../shared_preamble.tex (needs RedTitle, VeryLightGray)
%% and call \LecNineMap{k} inside a frame, where k in {1..5} highlights the
%% current act ("you are here") and k = 0 renders the closing reprise with all
%% five acts checked (used by T5's "Your Job Description" frame).
%% Filename deliberately does not match the build glob Lec*_T*.tex, so the
%% build script never tries to compile it standalone.

\newcommand{\LecNineMap}[1]{%
% Precompute per-box style selectors and the checkmark token; TikZ option
% lists are comma-split before TeX conditionals run, so \ifnum must not sit
% inside a [...] options list.
\ifnum#1=0
  \tikzset{lecmapa/.style={lecmapdone}, lecmapb/.style={lecmapdone},
           lecmapc/.style={lecmapdone}, lecmapd/.style={lecmapdone},
           lecmape/.style={lecmapdone}}%
  % green fill + the "all five acts complete" banner mark completion; a per-box
  % checkmark overflows the MANAGEMENT box, so keep the token empty here too.
  \def\lecmapchk{}%
\else
  \tikzset{lecmapa/.style={}, lecmapb/.style={}, lecmapc/.style={},
           lecmapd/.style={}, lecmape/.style={}}%
  \def\lecmapchk{}%
  \ifnum#1=1 \tikzset{lecmapa/.style={lecmapcur}}\fi
  \ifnum#1=2 \tikzset{lecmapb/.style={lecmapcur}}\fi
  \ifnum#1=3 \tikzset{lecmapc/.style={lecmapcur}}\fi
  \ifnum#1=4 \tikzset{lecmapd/.style={lecmapcur}}\fi
  \ifnum#1=5 \tikzset{lecmape/.style={lecmapcur}}\fi
\fi
\begin{center}
\begin{tikzpicture}[
    act/.style={rectangle, rounded corners, draw=black!60, fill=VeryLightGray,
        minimum width=2.6cm, minimum height=0.92cm, text width=2.5cm,
        align=center, font=\scriptsize, text=DarkText},
    lecmapcur/.style={draw=RedTitle, very thick, fill=orange!20},
    lecmapdone/.style={draw=green!45!black, thick, fill=green!10},
    q/.style={font=\tiny\itshape, text width=2.7cm, align=center, text=black!70},
    sat/.style={rectangle, rounded corners, draw=black!45, dashed, fill=white,
        text width=5.0cm, align=center, font=\tiny, text=black!70,
        minimum height=0.5cm},
    barlab/.style={font=\tiny\bfseries, text=black!60},
    arr/.style={-{Stealth[length=2mm]}, thick, gray!60}
]
    % --- five act boxes ---
    \node[act, lecmapa] (a1) at (0,0)
        {\textbf{9.1 MAP}\lecmapchk\\ \tiny what \& why};
    \node[act, lecmapb] (a2) at (2.85,0)
        {\textbf{9.2 MACHINE}\lecmapchk\\ \tiny under the hood};
    \node[act, lecmapc] (a3) at (5.7,0)
        {\textbf{9.3 METHOD}\lecmapchk\\ \tiny homotopy workflow};
    \node[act, lecmapd] (a4) at (8.55,0)
        {\textbf{9.4 MANAGEMENT}\lecmapchk\\ \tiny run the project};
    \node[act, lecmape] (a5) at (11.4,0)
        {\textbf{9.5 MINDSET}\lecmapchk\\ \tiny your role};
    \draw[arr] (a1) -- (a2);
    \draw[arr] (a2) -- (a3);
    \draw[arr] (a3) -- (a4);
    \draw[arr] (a4) -- (a5);
    % --- the student question each act answers ---
    \node[q] at (0,-0.82)    {what is this,\\ and why care?};
    \node[q] at (2.85,-0.82) {how does it work?\\ how do I start?};
    \node[q] at (5.7,-0.82)  {how do I structure\\ real work?};
    \node[q] at (8.55,-0.82) {how do I run a\\ whole project?};
    \node[q] at (11.4,-0.82) {what goes wrong?\\ what is my job?};
    % --- "you are here" marker above the current act ---
    \ifnum#1=1 \node[font=\scriptsize\bfseries, text=RedTitle] at (0,0.95)
        {you are here\ $\blacktriangledown$};\fi
    \ifnum#1=2 \node[font=\scriptsize\bfseries, text=RedTitle] at (2.85,0.95)
        {you are here\ $\blacktriangledown$};\fi
    \ifnum#1=3 \node[font=\scriptsize\bfseries, text=RedTitle] at (5.7,0.95)
        {you are here\ $\blacktriangledown$};\fi
    \ifnum#1=4 \node[font=\scriptsize\bfseries, text=RedTitle] at (8.55,0.95)
        {you are here\ $\blacktriangledown$};\fi
    \ifnum#1=5 \node[font=\scriptsize\bfseries, text=RedTitle] at (11.4,0.95)
        {you are here\ $\blacktriangledown$};\fi
    \ifnum#1=0 \node[font=\scriptsize\bfseries, text=green!45!black] at (5.7,0.95)
        {all five acts complete};\fi
    % --- bottom band: the three questions of the lecture ---
    \fill[blue!10]   (-1.3,-1.55) rectangle (4.15,-1.22);
    \fill[green!10]  (4.4,-1.55)  rectangle (9.85,-1.22);
    \fill[orange!15] (10.1,-1.55)  rectangle (12.7,-1.22);
    \node[barlab] at (1.425,-1.385)  {HOW IT WORKS};
    \node[barlab] at (7.125,-1.385)  {HOW TO USE IT WELL};
    \node[barlab] at (11.4,-1.385)   {YOUR ROLE};
    % --- dashed satellites: the two decks outside these five acts ---
    \node[sat] at (2.0,-2.2)
        {\textit{after 9.1:} optional interlude \textbf{9.1b} --- the agent landscape (MCP, A2A, benchmarks, safety)};
    \node[sat] at (9.8,-2.2)
        {\textit{after 9.5:} \textbf{9.6} wrap-up $\to$ \textbf{Lec10} full case studies (same morning)};
\end{tikzpicture}
\end{center}
}


\title[AI for Econ Research]{AI for Economic Research: Dynamic Models, Language, and Agents\\
Lecture 9.3: The Homotopy Workflow and Tool Choice}

\begin{document}

\makebeamertitle

%=============================================================================
\begin{frame}{The Map: Act 3 --- The Method}
\textbf{This deck answers the third question: how do I structure real work with the machine I just learned to drive?}

\LecNineMap{3}

\vspace{0.1cm}
\footnotesize\textit{Route: the homotopy idea $\to$ the six-step pipeline, step by step (with HA-Ramsey's actual opening prompts) $\to$ how to grade a plan before approving it $\to$ the runnable OLG demo $\to$ which tool for which step.}
\end{frame}

%=============================================================================
\section{The Homotopy Workflow}
%===========================================================

\sectiondivider{The Homotopy Workflow}{Build from simple to complex, validate, then extend}

\begin{frame}{The Homotopy Idea: Solve the Easy Version First}
\textbf{In workflow design, homotopy means solving an easier version first and then extending it toward the full research task.}

\vspace{0.15cm}

\begin{columns}[T]
\begin{column}{0.55\textwidth}
\begin{itemize}
    \item Start from a version with clear intuition and known checks
    \item Use that validated baseline as the starting point for the next
    \item Attribute new bugs to the new margin you just added
    \item Never ask the agent to solve the full complicated problem in one jump
\end{itemize}
\end{column}
\begin{column}{0.42\textwidth}
\begin{center}
\begin{tikzpicture}[
    v/.style={rectangle, draw=black, fill=VeryLightGray, rounded corners,
              minimum width=1.1cm, minimum height=0.55cm,
              text width=1.0cm, align=center, font=\scriptsize},
    chk/.style={font=\tiny, text=green!50!black},
    a/.style={-{Stealth[length=1.7mm]}, thick}
]
    \node[v] (v0) {V0\\det.};
    \node[v, right=0.45cm of v0] (v1) {V1\\+risk};
    \node[v, below=0.6cm of v0] (v2) {V2\\+dist.};
    \node[v, right=0.45cm of v2] (v3) {V3\\+fiscal};
    \draw[a] (v0) -- node[above,chk]{$\checkmark$} (v1);
    \draw[a] (v1) -- ++(0,-0.4) -- node[chk]{$\checkmark$} ++(-2.0,-0.2) -- (v2);
    \draw[a] (v2) -- node[above,chk]{$\checkmark$} (v3);
\end{tikzpicture}
\end{center}
\vspace{0.05cm}
\centering\scriptsize\textit{Each step small; each step validated; if a step fails, return to the last validated version.}
\end{column}
\end{columns}

\vspace{0.1cm}

\begin{conceptbox}
\footnotesize
\textbf{Why the name:} homotopy continuation starts from a problem with a known solution and \emph{continuously deforms} it toward the target (Judd 1998, Ch.~5--7). V0 $\to$ V1 $\to$ V2 are gates along that deformation --- the same discipline, scaled up, produced HA-Ramsey (Story 1; case in Lec10 T8).
\end{conceptbox}
\end{frame}

\begin{frame}{Homotopy in Practice: Build from Simple to Complex}
\begin{center}
\begin{tikzpicture}[
    box/.style={rectangle, draw=black, fill=VeryLightGray, rounded corners,
                minimum width=2.0cm, minimum height=0.95cm,
                text width=1.9cm, align=center, font=\small},
    arrow/.style={-{Stealth[length=2.5mm]}, thick}
]
    \node[box] (design) {1.\\Design};
    \node[box, right=0.25cm of design] (spec) {2.\\Write\\Spec};
    \node[box, right=0.25cm of spec] (algo) {3.\\Choose\\Algorithm};
    \node[box, right=0.25cm of algo] (v0) {4.\\Build\\V0};
    \node[box, right=0.25cm of v0] (val) {5.\\Validate};
    \node[box, right=0.25cm of val] (ext) {6.\\Extend};

    \draw[arrow] (design) -- (spec);
    \draw[arrow] (spec) -- (algo);
    \draw[arrow] (algo) -- (v0);
    \draw[arrow] (v0) -- (val);
    \draw[arrow] (val) -- (ext);
\end{tikzpicture}
\end{center}

\vspace{0.25cm}

\textbf{This replaces the fantasy prompt} ``solve my project.''

\begin{itemize}
    \item Start with a problem you can state and benchmark
    \item Give the agent bounded tasks with visible acceptance criteria
    \item Add one new ingredient at a time
\end{itemize}

\vspace{0.15cm}
\footnotesize\textit{The next frames walk the six steps in order --- with the HA-Ramsey case showing what each step looks like at research scale.}
\end{frame}

\begin{frame}{Step 1 --- Design: State the Goal Like an Economist}
\textbf{Before any code, write down the task in economist language.}

\vspace{0.2cm}

\begin{itemize}
    \item Research question in one sentence
    \item Objective and constraints, or the relevant source rules
    \item Equilibrium definition, estimand, or review standard if applicable
    \item Parameters, sample, source universe, and economic motivation
    \item Outputs you will treat as evidence
\end{itemize}

\vspace{0.15cm}

\textbf{Aiyagari anchor:}
\[
\max \; \mathbb{E}_0 \sum_{t=0}^{\infty} \beta^t u(c_t)
\quad \text{s.t.} \quad
c_t + a_{t+1} = (1+r)a_t + wz_t,\; a_{t+1} \geq \underline{a}
\]

\vspace{0.1cm}

\begin{takeawaybox}
\footnotesize
\textbf{Goal template} (T2's step 0, made precise): \emph{``Compute X with method Y to answer Z; success looks like W.''} If you cannot fill in W, you are not ready to open the terminal.
\end{takeawaybox}
\end{frame}

\begin{frame}{Step 2 --- The Spec Is the Contract}
\textbf{The specification file is the contract between you and the agent.}

\vspace{0.2cm}

\textbf{A good specification file states four things clearly:}
\begin{enumerate}
    \item assumptions and domain rules
    \item data or source restrictions
    \item procedure and validation target
    \item required outputs and where to save them
\end{enumerate}

\vspace{0.25cm}

\textbf{Examples:} CRRA utility and convergence tolerance for Aiyagari; weights and subgroup cuts for MEPS; quote-level evidence rules for paper review; official/institutional-source rules and provenance fields for fellows data.

\vspace{0.18cm}
\textbf{Best practice:} ask the agent to restate the specification in its own words before it writes code. That catches misunderstandings when they are still cheap.
\end{frame}

\begin{frame}[fragile]{How HA-Ramsey Opened: Restate, Audit, Refine}
\textbf{The restate-first practice is not hypothetical --- these three prompts opened the HA-Ramsey collaboration.}

\begin{lstlisting}[basicstyle=\ttfamily\scriptsize,numbers=none]
1. "Here are the RA paper (TeX) and the RA repository.
    Restate the model and the algorithm in your own words.
    List every equation and where it is implemented."
2. "Audit the code against the document. Report every place
    where the code and the TeX disagree: scoring, bounds,
    sampling, penalties. Classify each as a correctness risk
    or a documentation gap."
3. "Propose a refined RA document + code pair in which the
    two are consistent. Do not add any new economics yet."
\end{lstlisting}

\vspace{0.05cm}

\begin{itemize}
    \item Prompt 1 tests comprehension; prompt 2 tests the \emph{seed}; prompt 3 fixes the foundation while freezing the economics
    \item ``The audit is the cheapest insurance in the whole project'' --- extending an inconsistent seed propagates the inconsistency
\end{itemize}

\vspace{0.05cm}
{\scriptsize \textit{Source: HA-Ramsey case-study package (Lec10 T8).}}
\end{frame}

\begin{frame}{Grade the Plan Before Any Code: Five Criteria}
\textbf{T2 said: start in plan mode, review, approve. Approve against \emph{what}? These five criteria.}

\vspace{0.15cm}

\begin{enumerate}
    \item \textbf{Restates your goal correctly} --- in its own words, not your words echoed back
    \item \textbf{Bounded steps} --- each step has a named deliverable small enough to review in minutes
    \item \textbf{Acceptance criterion per step} --- a benchmark, test, or check named \emph{in advance}
    \item \textbf{Stop and rollback points} --- what happens when a step fails; which validated version you return to
    \item \textbf{Smallest testable object first} --- the plan begins with V0, not with the full target
\end{enumerate}

\vspace{0.15cm}

\begin{cautionbox}
\footnotesize
\textbf{A plan without acceptance criteria is a wish.} Send it back --- a plan revision costs one message; a wrong build costs a day.
\end{cautionbox}

\vspace{0.1cm}
\footnotesize\textit{Fellows instance: one batch of 50 = a bounded step; ``every filled field carries a source URL'' = its acceptance criterion. HA-Ramsey instance: the audit prompt is criterion 3 applied to the seed itself.}
\end{frame}

\begin{frame}[fragile]{Step 3 --- Algorithm and Pseudocode Catch Design Errors}
\textbf{Pseudocode or a phase plan is where you catch design mistakes before implementation.}

\vspace{0.2cm}

\begin{lstlisting}[style=pythonstyle,language={},basicstyle=\ttfamily\tiny,numbers=none]
Phase A: ingest or define the object
Phase B: run the chosen procedure
Phase C: save diagnostics and checks
Phase D: report only what passed validation
\end{lstlisting}

\vspace{0.15cm}

\textbf{Concrete procedures:} VFI or EGM for Aiyagari, a 3-prompt pipeline for MEPS, section decomposition for paper review, and batched metadata collection for fellows data.

\vspace{0.15cm}
\textbf{The point is not elegance.} The point is that you and the agent agree on the procedure before code details begin to hide conceptual errors.
\end{frame}

\begin{frame}{Step 4 --- V0: The Simplest Version With a Benchmark You Understand}
\textbf{V0 is the simplest version with a benchmark you understand --- ideally a closed form.}

\vspace{0.15cm}

\textbf{Optimal growth V0: full depreciation.} With log utility, Cobb--Douglas output $y_t = z_t k_t^{\alpha}$, and $\delta = 1$, the model has an exact closed-form solution:
\[
k_{t+1} = \alpha\beta\, z_t k_t^{\alpha}, \qquad c_t = (1-\alpha\beta)\, z_t k_t^{\alpha}
\]

\begin{itemize}
    \item Your V0 solver must reproduce this policy to numerical precision --- an \emph{absolute} benchmark, not a plausibility check
    \item Then relax $\delta < 1$: the closed form disappears, but the validated V0 has already certified the plumbing --- grids, interpolation, Euler equation, simulation
\end{itemize}

\vspace{0.15cm}

\textbf{Discipline:} do not ask the agent to add partial depreciation, risk, or equilibrium search before V0 matches the closed form.
\end{frame}

\begin{frame}{V0 at Research Scale: What Makes a Good Seed}
\textbf{At research scale, V0 has a name: the seed. HA-Ramsey grew from a published RA solver because that seed had four properties.}

\vspace{0.2cm}

\begin{enumerate}
    \item \textbf{It works} --- reproduces its own paper's numbers, not just ``runs without error''
    \item \textbf{It is documented} --- a TeX model note in the \emph{same notation} as the code
    \item \textbf{It is modular} --- config-driven, so every later experiment is a config diff, not a code fork
    \item \textbf{It points somewhere} --- the RA model is the degenerate no-risk case of the HA target; the deformation path exists by construction
\end{enumerate}

\vspace{0.2cm}

\begin{examplebox}
\footnotesize
\textbf{A good seed is like a good prototype for manufacturing a new product:} it already works end-to-end at small scale, so scaling up is engineering and refinement --- not a fresh invention on the factory floor.
\end{examplebox}

\vspace{0.1cm}
\textbf{Same object, different scales:} the classroom V0 and the research seed play the same role. The OLG demo's \texttt{v0} (later this deck) is your training-wheels seed.
\end{frame}

\begin{frame}{What ``V0 First'' Means Across Four Tasks}
\begin{center}
\small
\begin{tabular}{|p{2.2cm}|p{3.3cm}|p{5.2cm}|}
\hline
\textbf{Task} & \textbf{V0} & \textbf{First validation target} \\ \hline
Aiyagari model & deterministic household problem & recover simple savings logic and boundary behavior \\ \hline
MEPS pipeline & 3-year pilot panel & weighted rates line up with an external benchmark \\ \hline
Paper review skill & one section or theorem block & quotes are real and comments are specific \\ \hline
Fellows dataset & seed roster only & source rules and provenance fields are actually recorded \\ \hline
\end{tabular}
\end{center}

\vspace{0.25cm}
\textbf{Why economists should care:} V0 keeps the first error interpretable. You want the first failure to be obvious, not hidden inside a giant workflow.
\end{frame}

\begin{frame}{Steps 5--6 --- Validate, Then Extend One Margin at a Time}
\begin{center}
\small
\begin{tabular}{|c|p{3.5cm}|p{4.0cm}|}
\hline
\textbf{Version} & \textbf{New ingredient} & \textbf{Required benchmark} \\ \hline
V0 & Deterministic household & analytical cases \\ \hline
V1 & One new model, data, or review margin & setting the new margin to zero recovers V0 \\ \hline
V2 & Saved diagnostics and state tracking & mass sums, benchmarks, quote checks, or source logs behave sensibly \\ \hline
V3 & Full target workflow & final artifact survives validation and interpretation \\ \hline
\end{tabular}
\end{center}

\vspace{0.2cm}

\textbf{Why this works with agents:} when a new version breaks, both you and the agent know where to look.

\vspace{0.1cm}
\footnotesize\textit{Hold the gate idea: in T4 it returns as a management principle --- a quality gate on every intermediate good, not just on code versions.}
\end{frame}

\begin{frame}{Three Layers Grow Together: Model, Algorithm, Code}
\textbf{A research project is not one artifact but three, and the homotopy steps move all three in lockstep.}

\vspace{0.15cm}

\begin{center}
\begin{tikzpicture}[
    layer/.style={rectangle, rounded corners, draw=RedTitle!70, fill=blue!5,
        minimum width=3.1cm, minimum height=0.95cm, text width=3.0cm,
        align=center, font=\small},
    arr/.style={{Stealth[length=2mm]}-{Stealth[length=2mm]}, thick, gray!70},
    back/.style={-{Stealth[length=2mm]}, thick, dashed, green!45!black}
]
    \node[layer] (model) at (0,0) {\textbf{Model}\\ \scriptsize TeX: economics, timing, constraints};
    \node[layer] (algo) at (4.5,0) {\textbf{Algorithm}\\ \scriptsize pseudocode: solution method, penalties};
    \node[layer] (code) at (9.0,0) {\textbf{Code}\\ \scriptsize PyTorch/Python: what actually runs};
    \draw[arr] (model) -- (algo);
    \draw[arr] (algo) -- (code);
    \draw[back] (code.south) to[bend right=20] node[below, font=\tiny, text=green!45!black]
        {step back anytime: check, validate, revert} (model.south);
\end{tikzpicture}
\end{center}

\vspace{0.1cm}

\begin{itemize}
    \item \textbf{Forward:} you grow from the current version toward the final version gradually --- a new model element forces a new algorithm forces new code, one margin at a time
    \item \textbf{Backward:} every earlier version is preserved and validated, so you can always step back to a simpler one --- to check and validate a result, to isolate where something broke, or to revert to the last working version
    \item \textbf{The agent's job here:} keep TeX, pseudocode, and code describing \emph{the same economy} at every version --- drift between layers is where silent errors live
\end{itemize}

\vspace{0.05cm}
{\scriptsize \textit{Source: HA-Ramsey case-study package (Lec10 T8).}}
\end{frame}

\begin{frame}{Why This Beats One-Shot Generation}
\textbf{Four reasons the homotopy workflow is especially effective with agents:}

\begin{enumerate}
    \item \textbf{Agents excel at well-specified, bounded tasks.} \\
    ``Solve my model'' fails. ``Implement V0 per this pseudocode and validate against these 3 benchmarks'' succeeds.
    \medskip
    \item \textbf{Starting simple gives both you and the agent a working codebase.} \\
    V1 inherits validated V0 code. The agent builds on tested functions, not a blank file.
    \medskip
    \item \textbf{Bug attribution is immediate.} \\
    If V2 breaks, the bug is in V2's new lines, not the validated inherited ones.
    \medskip
    \item \textbf{Context builds naturally.} \\
    By V3, the agent has seen the spec, pseudocode, three validated versions, and every correction. \emph{Deep context} --- not a 10-page prompt.
\end{enumerate}

\textbf{Discipline:} resist the urge to skip ahead. A working V0 takes 30 minutes; debugging a broken V2 because V0 was never validated takes a day. This is also how good RAs are trained: simple task, verify, widen scope.
\end{frame}

%===========================================================


%=============================================================================
\section{OLG Five-Step Demo}
%===========================================================

\sectiondivider{OLG Five-Step Demo}{A concrete homotopy workflow for computational macro}

\begin{frame}{Demo: Folder Structure Is the Method}
\small
\begin{tabular}{p{4.0cm}p{7.7cm}}
\toprule
\textbf{Folder or file} & \textbf{Teaching role} \\
\midrule
\texttt{stages/} & five durable stages: model, equilibrium, algorithm, pseudocode, implementation/validation \\
\texttt{prompts/v*\_to\_v*.md} & one bounded terminal-agent session per version transition \\
\texttt{versions/v0--v5} & validated working code preserved at every step \\
\texttt{tests/} & one smoke test per version; failures identify the new margin \\
\texttt{build/*.json} & reproducibility reports from terminal runs \\
\bottomrule
\end{tabular}

\vspace{0.25cm}
\textbf{Demo path:} \texttt{demos/M4\_olg\_5step\_demo/} (local copy in this lecture; mirror of MiniCourse\_8hr \texttt{M4\_olg\_5step\_demo} minus \texttt{.venv}/\texttt{build/}).

\vspace{0.1cm}
\footnotesize\textit{Note the pattern: the HA-Ramsey \texttt{Case\_Study/} package (Lec10 T8) is this same folder discipline at paper scale --- seed, versions, notes, audits, all preserved.}
\end{frame}

\begin{frame}{Version Ladder: One Economic Margin at a Time}
\scriptsize
\begin{tabular}{p{0.8cm}p{5.0cm}p{5.8cm}}
\toprule
\textbf{Ver.} & \textbf{Adds} & \textbf{Validation gate} \\
\midrule
V0 & 3 cohorts, 1 asset, 2-state TFP, one NN policy & deterministic steady-state collapse; procyclical capital; RMS Euler residual below threshold \\
V1 & 7 cohorts and lifecycle labor & V0 logic preserved; middle-aged consumption peak appears \\
V2 & 4-state Markov TFP & collapse back to two states reproduces V1 \\
V3 & bonds, price, clearing, complementarity & zero bond-loss weight reproduces V2; positive late-life holdings \\
V4 & capital adjustment cost & zero adjustment cost reproduces V3; capital path smooths \\
V5 & stabilizing homotopy phases & residuals fall by phase; final Euler error target met \\
\bottomrule
\end{tabular}

\vspace{0.25cm}
\textbf{Research lesson:} the agent never receives ``solve the big model.'' It receives one validated deformation at a time.
\end{frame}

\begin{frame}[fragile]{Demo Launch and Classroom Contract}
\begin{lstlisting}[basicstyle=\ttfamily\scriptsize, numbers=none]
cd Lec09_Agentic_AI/demos/M4_olg_5step_demo
pip install -e .
python3 -m unittest discover -s tests
python3 run_demo.py
jupyter notebook demo.ipynb
\end{lstlisting}

\vspace{0.2cm}
\begin{itemize}
  \item \textbf{What the demo proves:} homotopy preserves a working baseline while complexity rises.
  \item \textbf{What to inspect:} prompts, version folders, test names, validation messages, and regenerated reports.
  \item \textbf{Failure recovery:} if V3 breaks, return to validated V2 and add only the bond-market-clearing layer.
\end{itemize}
\end{frame}


%=============================================================================
\section{Tool Choice}
%===========================================================

\sectiondivider{Tool Choice}{Match the tool to the pipeline step}

\begin{frame}{Tool Choice per Pipeline Step}
\textbf{The six-step pipeline also tells you which tool to open.}

\vspace{0.15cm}

\begin{center}
\small
\begin{tabular}{|p{2.4cm}|p{3.4cm}|p{5.2cm}|}
\hline
\textbf{Tool type} & \textbf{Pipeline steps} & \textbf{Economic example} \\ \hline
Web chat & 1--2: design, draft the spec & clarify identification logic, draft a model summary \\ \hline
IDE assistant & single-file edits inside a step & vectorize one function, explain one traceback \\ \hline
Terminal agent & 4--6: build, validate, extend & restructure the project, run scripts, generate diagnostics, iterate \\ \hline
\end{tabular}
\end{center}

\vspace{0.2cm}

\textbf{Why the terminal agent becomes the workhorse:} it operates on the real project and runs the real code --- while still leaving you in the approval and validation loop.

\vspace{0.1cm}
\textbf{For empirical research this is usually the sweet spot:} enough autonomy to save time, not so much autonomy that you stop checking the economics.
\end{frame}

\begin{frame}{Bridge: One Ladder vs.\ the Whole Factory}
\textbf{Act 3 disciplined the growth of one artifact. A real project is a factory: many artifacts, many sessions, several reviewers.}

\vspace{0.25cm}

\begin{itemize}
    \item You now have the \emph{method}: design, spec, algorithm, V0, validate, extend --- with a plan graded before any code
    \item But a paper is model + algorithm + code + data + reviews + write-up, worked across weeks of sessions
    \item The next question: \emph{``how do I run the whole project?''} --- Act 4's answer is a role change: the researcher as \textbf{project manager}, and the project objects (CLAUDE.md, skills, agents, rules) as the manager's instruments
\end{itemize}

\vspace{0.3cm}
\footnotesize\textit{Arc: T1 Framing $\to$ T2 Under the Hood $\to$ \textbf{T3 Homotopy} $\to$ T4 Project Org $\to$ T5 Mistakes/Context $\to$ T6 Wrap.}
\end{frame}

%===========================================================
\end{document}
